\documentclass[11pt]{article}
\usepackage{amsmath}
\usepackage{graphicx}
\usepackage{float}
\usepackage[backend=biber,sorting=none]{biblatex}
\usepackage{caption}
\usepackage[colorlinks=true,linkcolor=blue,urlcolor=black,bookmarksopen=true,citecolor=black]{hyperref}
\usepackage{hyperref}
\usepackage{cleveref}
\usepackage{multirow}
\usepackage{listings}
\usepackage{upquote}
\usepackage{color}
\usepackage{titling}
\usepackage{geometry}
 \geometry{
 a4paper,
 total={180mm,257mm},
 left=20mm,
 top=20mm,
 }

%\addbibresource{~/.config/LatexStuff/bibtex/bib/physics.bib}
\addbibresource{sources.bib}
\setlength{\parindent}{0pt}
\setlength{\droptitle}{-5em}

\title{Functional Programming: A1}
\author{Robert James Stevenson}
\date{08 May 2025}

\newcommand{\mycomment}[1]{}
\newcommand{\quotes}[1]{``#1''}
\definecolor{bluekeywords}{rgb}{0.13,0.13,1}
\definecolor{greencomments}{rgb}{0,0.5,0}
\definecolor{redstrings}{rgb}{0.9,0,0}

\lstdefinelanguage{FSharp}%
{morekeywords={let, new, match, with, rec, open, module, namespace, type, of, member, % 
and, for, while, true, false, in, do, begin, end, fun, function, return, yield, try, %
mutable, if, then, else, cloud, async, static, use, abstract, interface, inherit, finally },
  otherkeywords={ let!, return!, do!, yield!, use!, var, from, select, where, order, by },
  keywordstyle=\color{bluekeywords},
  sensitive=true,
  basicstyle=\footnotesize\ttfamily,
	breaklines=true,
  xleftmargin=\parindent,
  aboveskip=\bigskipamount,
	tabsize=4,
  morecomment=[l][\color{greencomments}]{///},
  morecomment=[l][\color{greencomments}]{//},
  morecomment=[s][\color{greencomments}]{{(*}{*)}},
  morestring=[b]",
  showstringspaces=false,
  literate={`}{\`}1,
  stringstyle=\color{redstrings},
}
\begin{document}
\maketitle
\section{Theory}
A high order function is a function whose arguments contain one or many functions. The function passed to the high order function is called a \quotes{callback}.
\subsection{Filter}
Filter iterates over every element in a list based data structure and returns a subset of the list, based on the conditions specified by the callback function.
At every iteration it calls the callback function (usually returns a boolean value) and passes the element to callback function. if the callback returns true,
the element is added to output list, else it's not added.
\\
Only the filter and list is passed to the function.
\subsection{Fold}
As the name suggests it folds (condenses) a collection (list) data structure into a single value(integer, string, another collection).
The reduction method is based on the callback function. It merges each element in collection into one value.
\\
Fold takes three arguments, they are the collection, the callback function and the accumulator (sometimes called initial value) which is returned by fold.
\subsection{Reduce}
Reduce is essentially is the same as Fold, except reduce doesn't have an accumulator argument, it assumes the first element in the collection as the 
accumulator and then iterates from the second value in the collection onwards (iteration starts at element 2). In my implementation reduce actually 
calls fold. \autocite{codecamp} \autocite{Philip} \autocite{geeks}.
\section{Implementation}
\subsection{Map}
\lstinputlisting[language=FSharp,linerange={1-6}]{A1.fsx}
\subsection{Filter}
\lstinputlisting[language=FSharp,linerange={8-16}]{A1.fsx}
\subsection{Fold}
\lstinputlisting[language=FSharp,linerange={19-25}]{A1.fsx}
\subsection{Reduce}
\lstinputlisting[language=FSharp,linerange={28-44}]{A1.fsx}

\printbibliography
\end{document}

